\chapter{Entering a New Market}\label{ch:fm:entering-a-new-market}

The plan says eight months to the first trade. The clearing agreement alone takes five, and it cannot start until the legal entity exists and
has its identifier; the exchange will not certify the software until the clearing firm has set the firm's limits; and the regulatory registration
that membership requires takes four months of its own, starting only when the entity exists. Drawn as a network, the entry takes ten months, and the
long pole is not the clearing agreement everyone was watching but the registration nobody had put on the chart.

\section{Why enter: the edge and the size of the prize}

A trading firm enters a new market (an asset class, a region, a venue) because an edge it has elsewhere should work there, or because a client or
counterparty it already serves trades there. Flow Traders' 2025 annual report is a public example of the second kind of reasoning: it described
expanding its presence in Asia and ``initiating active trading in China, an attractive growth market aligned with our core ETP strengths'', with its
exchange-traded product value traded in Asia rising from \euro114 billion in 2024 to \euro152 billion in 2025. Whatever the reason, the entry costs
money for months before the first trade, and the prize is uncertain until the firm has traded for a while.

\begin{definition}[Entry plan, time to first trade]\label{def:fm:entering-a-new-market:plan}
An \emph{entry plan}\index{entry plan} sets out the workstreams needed to trade in a new market (legal entity, \omterm{def:fm:regulators-and-licences:perimeter}{authorisation}, clearing,
membership, connectivity, data, people, software, certification), their durations, costs and dependencies, the expected revenue ramp and the
points at which the firm decides whether to continue. The \emph{time to first trade}\index{time to first trade} is the time from the decision to
enter to the first live trade, set by the longest chain of dependent workstreams.
\end{definition}

\section{The checklist: access, clearing, data, legal, staffing}

The workstreams are the chapters of this series. The entity needs a legal entity identifier (Book 2) before most counterparties will open an
account; \omterm{def:fm:regulators-and-licences:perimeter}{authorisation} or registration follows chapter 17's map and its statutory periods; clearing through a futures commission merchant or a
clearing member (Book 1) needs know-your-customer checks, legal agreements (chapter 18) and limits; membership and certification (Book 13) follow;
connectivity (Book 14), market data licences (chapter 22) and people (chapter 10) run alongside. The chapter's example is a futures market in a new
region (\cref{tab:fm:entering-a-new-market:tasks}); durations and costs are illustrative inputs.

\begin{table}[ht]
\centering\small
\begin{tabular}{lrrlr}
\toprule
workstream & months & \$k/month & after & slack \\
\midrule
legal entity & 2.0 & 25 & -- & 0 \\
legal entity identifier & 0.5 & 2 & entity & 1.0 \\
regulatory registration & 4.0 & 30 & entity & 0 \\
clearing agreement & 5.0 & 20 & identifier & 1.0 \\
clearing limits set & 0.5 & 5 & clearing agreement & 1.0 \\
exchange membership & 3.0 & 15 & registration & 0 \\
connectivity & 3.0 & 40 & entity & 4.0 \\
market data licences & 2.0 & 10 & entity & 6.0 \\
hires & 4.0 & 60 & -- & 2.0 \\
software adaptation & 3.0 & 80 & hires & 2.0 \\
certification & 1.0 & 10 & limits, membership, & 0 \\
 & & & connectivity, software & \\
\bottomrule
\end{tabular}
\caption{The entry's workstreams, with their slack in months under the earliest schedule (illustrative). Data: \texttt{fm\_entry.TASKS}, \texttt{firm.entryplan.schedule}.}\label{tab:fm:entering-a-new-market:tasks}
\end{table}

\section{The critical path}

\begin{definition}[Critical-path method]\label{def:fm:entering-a-new-market:cpm}
The \emph{critical-path method}\index{critical-path method} schedules a project given as tasks with durations and precedence constraints: each task
starts when all its predecessors finish; the project's length is the longest path through the network; a task's slack is how much it can be delayed
without delaying the end; the critical path is the chain of tasks with no slack.
\end{definition}

Kelley and Walker described the method in 1959 for construction and maintenance projects; it applies unchanged to an \omterm{def:fm:entering-a-new-market:plan}{entry plan}
(\cref{lst:fm:entering-a-new-market:cpm}).

\omcode{code/firm/entryplan/firm_entryplan.py}{30}{64}{The critical-path method: earliest starts and finishes by recursion, latest finishes backwards,
slack, and the critical path; and the monthly cost schedule.}\label{lst:fm:entering-a-new-market:cpm}

The critical path runs through the legal entity, the regulatory registration, the exchange membership and the certification: ten months to the first
trade (\cref{fig:fm:entering-a-new-market:gantt}). The clearing chain, which the plan's eight months were built around, finishes after eight
months and has a month of slack. Connectivity has four months of slack and the data licences six: they can start late without cost. Before the first
trade the entry spends \$949\,000, at most \$170\,000 in a month, when software, hires, registration and connectivity overlap.

\begin{omfigure}
\begin{tikzpicture}
\begin{axis}[width=10.5cm,height=6.2cm,xbar stacked,bar width=7pt,xmin=0,xmax=11,xlabel={\small month},
  ytick={0,1,2,3,4,5,6,7,8,9,10},yticklabels={legal entity,identifier,registration,clearing agreement,clearing limits,membership,connectivity,
  data licences,hires,software,certification},y dir=reverse,ymin=-0.6,ymax=10.6,tick label style={font=\footnotesize},grid=major,
  grid style={gray!20},table/col sep=comma]
\addplot[draw=none,fill=none] table[y=pos,x=start,restrict expr to domain={\coordindex}{0:10}]{figdata/desk/24-entering-a-new-market/gantt.csv};
\addplot[fill=omThm!45,draw=omThm] table[y=pos,x expr=\thisrow{months}*(1-\thisrow{critical}),restrict expr to domain={\coordindex}{0:10}]{figdata/desk/24-entering-a-new-market/gantt.csv};
\addplot[fill=omDef!70,draw=omDef] table[y=pos,x expr=\thisrow{months}*\thisrow{critical},restrict expr to domain={\coordindex}{0:10}]{figdata/desk/24-entering-a-new-market/gantt.csv};
\draw[black,thick] (axis cs:10,-0.6) -- (axis cs:10,10.6) node[pos=0.03,left,font=\scriptsize] {first trade};
\draw[gray,dashed] (axis cs:8,-0.6) -- (axis cs:8,10.6) node[pos=0.97,left,font=\scriptsize] {plan};
\end{axis}
\end{tikzpicture}

\omcaption{The entry's schedule: each workstream from its earliest start, the critical path in the darker colour. The first trade comes at ten
months, not the plan's eight: the registration and membership chain is two months longer than the clearing chain. Data:
\texttt{firm.entryplan.schedule}.}\label{fig:fm:entering-a-new-market:gantt}
\end{omfigure}

\section{Staging the investment}

\begin{definition}[Staged investment]\label{def:fm:entering-a-new-market:staged}
A \emph{staged investment}\index{staged investment} commits money in steps separated by decision points, each with a kill criterion (Book 7's
stage gates) stated in advance, so that the firm can stop when what it has learnt says the rest is not worth spending.
\end{definition}

After the first trade the revenue ramps towards a steady state $R$ a month, as $R(1-e^{-t/\tau})$ with $\tau=4$ months, against running costs of
\$250\,000 a month; the horizon is three years and the discount rate 10\%. The steady state is the uncertain part: in the chapter's model it is
lognormal with a median of \$400\,000 a month. A weak outcome (\$150\,000) loses \$3.55 million of present value; the base case (\$450\,000) gains
\$2.00 million, an internal rate of return of 104\% a year; a strong one (\$800\,000) gains \$8.46 million.

\begin{proposition}[A kill point never hurts with perfect information]\label{prop:fm:entering-a-new-market:kill}
If at the kill point the firm learns $R$ exactly and stops when $R$ is below the running cost $c$, the entry's expected net present value with the
kill point is at least its value without.
\end{proposition}

\begin{proof}
After the kill point, continuing adds $R(1-e^{-t/\tau})-c<R-c<0$ each month when $R<c$, and stopping adds zero; when $R\ge c$ the firm continues and
nothing changes. The kill point therefore raises the value in every state with $R<c$ and leaves the others unchanged; so does its expectation.
\end{proof}

In practice the firm sees a noisy estimate. The chapter's firm observes $R$ with a 25\% error after six months of live trading and stops if the
estimate is below the \$250\,000 running cost (\cref{lst:fm:entering-a-new-market:stage}).

\omcode{code/firm/entryplan/firm_entryplan.py}{94}{107}{The value of staging: for each true revenue, the NPV without a kill point and with one based on
a noisy estimate.}\label{lst:fm:entering-a-new-market:stage}

\section{Tutorial: eight months to first trade}\label{tut:fm:entering-a-new-market:tut}

\begin{tutorial}
\textbf{Goal.} Draw the entry's critical path, price it, and value a kill point after six months of trading. \textbf{End state:} the Gantt chart
(\cref{fig:fm:entering-a-new-market:gantt}) and the NPV distributions with and without staging (\cref{fig:fm:entering-a-new-market:npv}).
\begin{tutsteps}
\item \textbf{The network.} \texttt{fm\_entry.TASKS}; \texttt{firm.entryplan.schedule} and \texttt{critical\_path}.
\item \textbf{The costs.} \texttt{cost\_by\_month} spreads each workstream's cost over its months.
\item \textbf{The outcomes.} \texttt{fm\_entry.scenario\_npvs()} for three revenue levels; \texttt{revenue\_draws} for 4\,000 lognormal draws.
\item \textbf{The kill point.} \texttt{fm\_entry.staging()} compares the NPVs; \texttt{threshold\_curve()} varies the kill threshold.
\end{tutsteps}
\end{tutorial}

\begin{omfigure}
\begin{tikzpicture}
\begin{axis}[width=10.5cm,height=5.2cm,ybar,bar width=4pt,xlabel={\small net present value (\$ million)},ylabel={\small share of outcomes (\%)},
  xmin=-6.5,xmax=15,ymin=0,ymax=30,tick label style={font=\footnotesize},grid=major,grid style={gray!20},table/col sep=comma,
  legend cell align=left,legend style={font=\scriptsize,at={(0.5,-0.3)},anchor=north,/tikz/every even column/.append style={column sep=8pt}},legend columns=2]
\addplot[fill=gray!45,draw=gray] table[x=mid,y=plain]{figdata/desk/24-entering-a-new-market/npv.csv};
\addplot[fill=omDef!55,draw=omDef] table[x=mid,y=staged]{figdata/desk/24-entering-a-new-market/npv.csv};
\legend{no kill point,kill point after six months}
\end{axis}
\end{tikzpicture}

\omcaption{The entry's net present value over 4\,000 draws of its steady-state revenue, with and without a kill point after six months of trading.
Staging removes the worst losses: the 5th percentile rises from $-\$3.57$ million to $-\$1.92$ million. Data: \texttt{fm\_entry.staging}.}\label{fig:fm:entering-a-new-market:npv}
\end{omfigure}

Without a kill point the entry's expected NPV is \$2.52 million, with a 38.6\% chance of losing money and a 5th percentile of $-\$3.57$ million. With
the kill point the firm stops 22.7\% of entries; the expected NPV rises by \$200\,000 to \$2.72 million, and the 5th percentile to $-\$1.92$
million. The probability of losing money barely moves (39.6\%): stopping turns large losses into smaller ones rather than into gains, and the
set-up costs are sunk. The noise costs something: 3.6\% of all entries are weak and continue, and 5.5\% are sound and stopped.

\begin{omfigure}
\begin{tikzpicture}
\begin{axis}[width=10.5cm,height=5cm,xlabel={\small kill threshold (\$ thousand of estimated monthly revenue)},ylabel={\small value of the kill point (\$ thousand)},
  xmin=0,xmax=700,ymin=-1400,ymax=300,tick label style={font=\footnotesize},grid=major,grid style={gray!20},table/col sep=comma]
\addplot[omDef,thick,mark=*,mark size=1.2pt] table[x=threshold,y=value]{figdata/desk/24-entering-a-new-market/threshold.csv};
\draw[black,thin] (axis cs:0,0) -- (axis cs:700,0);
\draw[gray,dotted] (axis cs:250,-1400) -- (axis cs:250,300) node[pos=0.1,right,font=\scriptsize] {running cost};
\end{axis}
\end{tikzpicture}

\omcaption{The value the kill point adds to the entry's expected NPV, by the revenue threshold below which the firm stops. It peaks near the
running cost, \$250\,000 a month, and turns negative above \$350\,000, where the rule stops too many sound entries. Data:
\texttt{fm\_entry.threshold\_curve}.}\label{fig:fm:entering-a-new-market:threshold}
\end{omfigure}

The threshold matters as much as the kill point (\cref{fig:fm:entering-a-new-market:threshold}). Set at the running cost, as
\cref{prop:fm:entering-a-new-market:kill} suggests, it adds \$200\,000; at \$300\,000 it adds \$150\,000; at \$400\,000 it destroys \$138\,000,
because a rule that demands the base case kills entries that would have paid. A kill criterion is written before entry, at a level the firm could
defend if the market then turned out well.

\section{A worked entry plan}

\begin{method}[Planning an entry]\label{met:fm:entering-a-new-market:plan}
\begin{enumerate}
\item State why the firm's edge should work in the market, and the size of the prize, with evidence (Book 7's research standards).
\item List every workstream with its owner, duration, cost and predecessors; draw the network and its critical path; put the owners of the critical
workstreams in the weekly meeting.
\item Start the long, cheap workstreams first (registration, clearing); start the short, expensive ones (hires, connectivity) as late as their slack
allows.
\item Write the kill criteria before entry: when they will be measured, on what data, and the threshold.
\item Review the plan at every stage gate against the network and the criteria, not against the original dates.
\end{enumerate}
\end{method}

\begin{remark}[Which slack to use]\label{rem:fm:entering-a-new-market:slack}
Slack is an option on delay. Starting the hires two months later, within their slack, defers \$240\,000 of salaries without delaying the first
trade, and keeps the option not to spend it if the registration or the clearing agreement fails first. The same logic applies to connectivity,
with four months of slack and fees that start on delivery: the cheapest workstream to start is the one that decides the date.
\end{remark}

\section{Build: the entry plan}\label{bld:fm:entering-a-new-market:entryplan}

\begin{build}
\bfield{Purpose} An entry as a network of workstreams: the critical path and slack, the cost schedule, revenue ramps, NPV and IRR, and the value of
staging.
\bfield{Interface} \texttt{firm.entryplan}: \texttt{Task}, \texttt{schedule}, \texttt{critical\_path}, \texttt{cost\_by\_month}, \texttt{npv},
\texttt{irr}, \texttt{entry\_cash}, \texttt{staged\_npv}.
\bfield{Rules} A task starts when all its predecessors finish; the critical path has zero slack; a killed entry stops its running costs and revenue
at the kill month.
\bfield{Acceptance tests} \texttt{code/firm/entryplan/tests/}: a four-task network's schedule, slack and path; costs, NPV and IRR on hand numbers;
the kill rule with exact information.
\bfield{Stretch} Uncertain durations and the probability of meeting a date; costs from Book 1's clearing and fee models; several kill points.
\end{build}

\begin{omsources}
\item Flow Traders, Annual Report 2025.
\item J.~E.~Kelley and M.~R.~Walker, ``Critical-path planning and scheduling'', Eastern Joint Computer Conference, 1959.
\item A.~K.~Dixit and R.~S.~Pindyck, \emph{Investment under Uncertainty}, 1994.
\end{omsources}

\section{Exercises}

\begin{exercise}[$\star$]\label{exo:fm:entering-a-new-market:1}
Compute the length of the clearing chain and of the registration chain. Which is critical?
\end{exercise}

\begin{exercise}[$\star$]\label{exo:fm:entering-a-new-market:2}
Which workstreams have the most slack, and what does that allow?
\end{exercise}

\begin{exercise}[$\star$]\label{exo:fm:entering-a-new-market:3}
What does the entry cost before the first trade, and in which month does it spend most?
\end{exercise}

\begin{exercise}[$\star\star$]\label{exo:fm:entering-a-new-market:4}
If the registration took two months instead of four, what would the \omterm{def:fm:entering-a-new-market:plan}{time to first trade} be, and which chain would be critical?
\end{exercise}

\begin{exercise}[$\star\star$]\label{exo:fm:entering-a-new-market:5}
Prove \cref{prop:fm:entering-a-new-market:kill}, and explain why it can fail with a noisy estimate.
\end{exercise}

\begin{exercise}[$\star\star$]\label{exo:fm:entering-a-new-market:6}
Why does the kill point barely change the probability of losing money while raising the 5th percentile by \$1.6 million?
\end{exercise}

\begin{exercise}[$\star\star\star$]\label{exo:fm:entering-a-new-market:7}
\emph{Coding.} Rerun \texttt{fm\_entry.threshold\_curve()} and find the best threshold. What does a threshold of \$400\,000 do?
\end{exercise}

\begin{exercise}[$\star\star\star$]\label{exo:fm:entering-a-new-market:8}
\emph{Find the flaw.} ``The clearing agreement is the long pole: we will start it first and the rest can wait.''
\end{exercise}

\section{Problem: Eight Months to First Trade}

\begin{problem}[{Weekend problem --- eight months to first trade}]\label{pb:fm:entering-a-new-market:1}
A head of business development must present an \omterm{def:fm:entering-a-new-market:plan}{entry plan} for a futures market in a new region, with a date and a decision rule.

\medskip\noindent\textbf{Part I --- Why and what.}
\begin{enumerate}
\item Why would a trading firm enter a new market? Give the public example of the chapter.
\item Define an \omterm{def:fm:entering-a-new-market:plan}{entry plan} and the \omterm{def:fm:entering-a-new-market:plan}{time to first trade}.
\item List the workstreams and the chapter each draws on.
\item Which dependencies does the hook describe?
\end{enumerate}

\medskip\noindent\textbf{Part II --- The network.}
\begin{enumerate}[resume]
\item Define the \omterm{def:fm:entering-a-new-market:cpm}{critical-path method}.
\item Give the schedule and the critical path.
\item Give the slack of each workstream and what to do with it.
\item Give the cost before the first trade and its peak month.
\end{enumerate}

\medskip\noindent\textbf{Part III --- The money.}
\begin{enumerate}[resume]
\item Describe the revenue ramp and the three scenarios' NPVs.
\item Define a \omterm{def:fm:entering-a-new-market:staged}{staged investment}.
\item State and prove \cref{prop:fm:entering-a-new-market:kill}.
\item Give the NPV distribution with and without the kill point.
\item How should the kill threshold be set?
\end{enumerate}

\medskip\noindent\textbf{Part IV --- The plan.}
\begin{enumerate}[resume]
\item Which workstreams would you start first, and which last?
\item Who should attend the weekly meeting?
\item What would you report at each stage gate?
\item What would make you shorten the kill point to three months?
\item What does the model leave out?
\item State the \emph{named result}: the \omterm{def:fm:entering-a-new-market:plan}{time to first trade} on the critical path, and the value that the kill point adds to the entry's NPV.
\item In two sentences, write the plan.
\end{enumerate}
\end{problem}

\section{Interview questions}

\begin{interviewq}[$\star$ \iqroles{developer}]\label{iq:fm:entering-a-new-market:1}
Compute the critical path of a small project from its tasks and dependencies.
\end{interviewq}

\begin{interviewq}[$\star$ \iqroles{trader}]\label{iq:fm:entering-a-new-market:2}
What has to be in place before a firm can trade futures on a new exchange?
\end{interviewq}

\begin{interviewq}[$\star\star$ \iqroles{researcher}]\label{iq:fm:entering-a-new-market:3}
How would you value the option to stop a new business after six months?
\end{interviewq}

\begin{interviewq}[$\star\star$ \iqroles{risk}]\label{iq:fm:entering-a-new-market:4}
What kill criteria would you set for a new market entry, and why before entering?
\end{interviewq}

\begin{interviewq}[$\star\star$ \iqroles{trader, researcher}]\label{iq:fm:entering-a-new-market:5}
Why might a firm that is profitable elsewhere lose money entering a new region?
\end{interviewq}

\begin{interviewq}[$\star\star\star$ \iqroles{developer, researcher}]\label{iq:fm:entering-a-new-market:6}
Durations are uncertain. How would you estimate the probability of trading by a given date?
\end{interviewq}
